\chapter{Considerações finais}

Este trabalho teve como foco caracterizar a adequação prática e os limites de um \textit{pipeline} avaliativo de reconstrução de cenários 3D fotorrealistas para VR \textit{standalone}, articulando modelos NeRF e \textit{3D Gaussian Splatting}. A motivação central partiu do fato de que técnicas recentes de renderização neural alcançam elevada fidelidade visual, mas ainda encontram limitações importantes quando precisam ser integradas a uma experiência imersiva, estereoscópica e executada sob restrições de hardware móvel, como ocorre no Meta Quest 3S.

Ao longo do trabalho, a fundamentação teórica permitiu situar o problema no contexto da síntese de novas vistas, da reconstrução multivista e das diferenças entre representações implícitas e explícitas. O NeRF foi discutido como linha de base de fidelidade visual, enquanto o 3D-GS foi apresentado como alternativa mais próxima dos requisitos de renderização interativa. Essa comparação evidenciou que o desafio do projeto não está em criar uma nova técnica de 3D-GS, mas em adaptar, integrar e avaliar a representação em um cenário de VR \textit{standalone}.

Com base nessa análise, foi proposto um projeto organizado em etapas, partindo da captura de imagens RGB e do pré-processamento dos dados, avançando para o treinamento de um modelo NeRF, a reconstrução ou conversão em 3D-GS, a integração em Unreal Engine e a avaliação em ambiente imersivo. A definição de requisitos, arquitetura, fluxo de implementação, métricas e riscos contribuiu para transformar o problema em um plano experimental controlável, no qual qualidade visual, FPS, tempo de quadro, uso de memória, estabilidade, artefatos e estranhamento perceptual são analisados como dimensões complementares de adequação ao cenário-alvo.

Também ficaram delimitadas as principais restrições do estudo. O trabalho concentra-se em cenas estáticas, capturadas por imagens RGB, e não contempla reconstrução dinâmica, \textit{relighting} avançado ou interação física complexa com os elementos da cena. Além disso, a entrega prevista é um protótipo avaliativo, e não uma ferramenta completa de captura, edição e publicação de cenas 3D-GS. Possíveis ferramentas integradas ficam posicionadas como trabalhos futuros.

Dessa forma, esta etapa do trabalho consolida a base conceitual e metodológica necessária para a continuidade no TCC III. Como continuidade, espera-se implementar e validar o \textit{pipeline} proposto, coletar métricas de qualidade, desempenho e estabilidade perceptual, analisar as limitações encontradas e indicar aprimoramentos para futuras aplicações de reconstrução neural em VR \textit{standalone}.
